% Options for packages loaded elsewhere
% Options for packages loaded elsewhere
\PassOptionsToPackage{unicode}{hyperref}
\PassOptionsToPackage{hyphens}{url}
\PassOptionsToPackage{dvipsnames,svgnames,x11names}{xcolor}
%
\documentclass[
  12pt,
  letterpaper,
  DIV=11,
  numbers=noendperiod]{scrartcl}
\usepackage{xcolor}
\usepackage{amsmath,amssymb}
\setcounter{secnumdepth}{5}
\usepackage{iftex}
\ifPDFTeX
  \usepackage[T1]{fontenc}
  \usepackage[utf8]{inputenc}
  \usepackage{textcomp} % provide euro and other symbols
\else % if luatex or xetex
  \usepackage{unicode-math} % this also loads fontspec
  \defaultfontfeatures{Scale=MatchLowercase}
  \defaultfontfeatures[\rmfamily]{Ligatures=TeX,Scale=1}
\fi
\usepackage{lmodern}
\ifPDFTeX\else
  % xetex/luatex font selection
  \setmainfont[]{Times New Roman}
\fi
% Use upquote if available, for straight quotes in verbatim environments
\IfFileExists{upquote.sty}{\usepackage{upquote}}{}
\IfFileExists{microtype.sty}{% use microtype if available
  \usepackage[]{microtype}
  \UseMicrotypeSet[protrusion]{basicmath} % disable protrusion for tt fonts
}{}
\makeatletter
\@ifundefined{KOMAClassName}{% if non-KOMA class
  \IfFileExists{parskip.sty}{%
    \usepackage{parskip}
  }{% else
    \setlength{\parindent}{0pt}
    \setlength{\parskip}{6pt plus 2pt minus 1pt}}
}{% if KOMA class
  \KOMAoptions{parskip=half}}
\makeatother
% Make \paragraph and \subparagraph free-standing
\makeatletter
\ifx\paragraph\undefined\else
  \let\oldparagraph\paragraph
  \renewcommand{\paragraph}{
    \@ifstar
      \xxxParagraphStar
      \xxxParagraphNoStar
  }
  \newcommand{\xxxParagraphStar}[1]{\oldparagraph*{#1}\mbox{}}
  \newcommand{\xxxParagraphNoStar}[1]{\oldparagraph{#1}\mbox{}}
\fi
\ifx\subparagraph\undefined\else
  \let\oldsubparagraph\subparagraph
  \renewcommand{\subparagraph}{
    \@ifstar
      \xxxSubParagraphStar
      \xxxSubParagraphNoStar
  }
  \newcommand{\xxxSubParagraphStar}[1]{\oldsubparagraph*{#1}\mbox{}}
  \newcommand{\xxxSubParagraphNoStar}[1]{\oldsubparagraph{#1}\mbox{}}
\fi
\makeatother


\usepackage{longtable,booktabs,array}
\newcounter{none} % for unnumbered tables
\usepackage{calc} % for calculating minipage widths
% Correct order of tables after \paragraph or \subparagraph
\usepackage{etoolbox}
\makeatletter
\patchcmd\longtable{\par}{\if@noskipsec\mbox{}\fi\par}{}{}
\makeatother
% Allow footnotes in longtable head/foot
\IfFileExists{footnotehyper.sty}{\usepackage{footnotehyper}}{\usepackage{footnote}}
\makesavenoteenv{longtable}
\usepackage{graphicx}
\makeatletter
\newsavebox\pandoc@box
\newcommand*\pandocbounded[1]{% scales image to fit in text height/width
  \sbox\pandoc@box{#1}%
  \Gscale@div\@tempa{\textheight}{\dimexpr\ht\pandoc@box+\dp\pandoc@box\relax}%
  \Gscale@div\@tempb{\linewidth}{\wd\pandoc@box}%
  \ifdim\@tempb\p@<\@tempa\p@\let\@tempa\@tempb\fi% select the smaller of both
  \ifdim\@tempa\p@<\p@\scalebox{\@tempa}{\usebox\pandoc@box}%
  \else\usebox{\pandoc@box}%
  \fi%
}
% Set default figure placement to htbp
\def\fps@figure{htbp}
\makeatother





\setlength{\emergencystretch}{3em} % prevent overfull lines

\providecommand{\tightlist}{%
  \setlength{\itemsep}{0pt}\setlength{\parskip}{0pt}}



 


\usepackage{tcolorbox}
\usepackage{amssymb}
\usepackage{yfonts}
\usepackage{bm}


\newtcolorbox{greybox}{
  colback=white,
  colframe=blue,
  coltext=black,
  boxsep=5pt,
  arc=4pt}
  
\newcommand{\sectionbreak}{\clearpage}

 
\newcommand{\ds}[4]{\sum_{{#1}=1}^{#3}\sum_{{#2}=1}^{#4}}
\newcommand{\us}[3]{\mathop{\sum\sum}_{1\leq{#2}<{#1}\leq{#3}}}

\newcommand{\ol}[1]{\overline{#1}}
\newcommand{\ul}[1]{\underline{#1}}

\newcommand{\amin}[1]{\mathop{\text{argmin}}_{#1}}
\newcommand{\amax}[1]{\mathop{\text{argmax}}_{#1}}

\newcommand{\ci}{\perp\!\!\!\perp}

\newcommand{\mc}[1]{\mathcal{#1}}
\newcommand{\mb}[1]{\mathbb{#1}}
\newcommand{\mf}[1]{\mathfrak{#1}}

\newcommand{\eps}{\epsilon}
\newcommand{\lbd}{\lambda}
\newcommand{\alp}{\alpha}
\newcommand{\df}{=:}
\newcommand{\am}[1]{\mathop{\text{argmin}}_{#1}}
\newcommand{\ls}[2]{\mathop{\sum\sum}_{#1}^{#2}}
\newcommand{\ijs}{\mathop{\sum\sum}_{1\leq i<j\leq n}}
\newcommand{\jis}{\mathop{\sum\sum}_{1\leq j<i\leq n}}
\newcommand{\sij}{\sum_{i=1}^n\sum_{j=1}^n}
	
\KOMAoption{captions}{tableheading}
\makeatletter
\@ifpackageloaded{caption}{}{\usepackage{caption}}
\AtBeginDocument{%
\ifdefined\contentsname
  \renewcommand*\contentsname{Table of contents}
\else
  \newcommand\contentsname{Table of contents}
\fi
\ifdefined\listfigurename
  \renewcommand*\listfigurename{List of Figures}
\else
  \newcommand\listfigurename{List of Figures}
\fi
\ifdefined\listtablename
  \renewcommand*\listtablename{List of Tables}
\else
  \newcommand\listtablename{List of Tables}
\fi
\ifdefined\figurename
  \renewcommand*\figurename{Figure}
\else
  \newcommand\figurename{Figure}
\fi
\ifdefined\tablename
  \renewcommand*\tablename{Table}
\else
  \newcommand\tablename{Table}
\fi
}
\@ifpackageloaded{float}{}{\usepackage{float}}
\floatstyle{ruled}
\@ifundefined{c@chapter}{\newfloat{codelisting}{h}{lop}}{\newfloat{codelisting}{h}{lop}[chapter]}
\floatname{codelisting}{Listing}
\newcommand*\listoflistings{\listof{codelisting}{List of Listings}}
\makeatother
\makeatletter
\makeatother
\makeatletter
\@ifpackageloaded{caption}{}{\usepackage{caption}}
\@ifpackageloaded{subcaption}{}{\usepackage{subcaption}}
\makeatother
\usepackage{bookmark}
\IfFileExists{xurl.sty}{\usepackage{xurl}}{} % add URL line breaks if available
\urlstyle{same}
\hypersetup{
  pdftitle={Array Decomposition},
  pdfauthor={Jan de Leeuw},
  colorlinks=true,
  linkcolor={blue},
  filecolor={Maroon},
  citecolor={Blue},
  urlcolor={Blue},
  pdfcreator={LaTeX via pandoc}}


\title{Array Decomposition}
\author{Jan de Leeuw}
\date{May 19, 2026}
\begin{document}
\maketitle
\begin{abstract}
TBD
\end{abstract}

\renewcommand*\contentsname{Table of contents}
{
\setcounter{tocdepth}{3}
\tableofcontents
}

\sectionbreak

\textbf{Note:} This is a working manuscript which will be
expanded/updated frequently. All suggestions for improvement are
welcome. All Rmd, tex, html, pdf, R, and C files are in the public
domain. Attribution will be appreciated, but is not required. The files
can be found at \url{https://github.com/deleeuw/arrays}

\sectionbreak

\section{Introduction}\label{introduction}

\section{Decomposition}\label{decomposition}

\subsection{Vectors}\label{vectors}

Suppose \(x\) is an element of \(n\)-dimensional Euclidean space, with
the usual inner product and norm. We use a \(\bullet\) replacing a
subscript to indicate averaging over that subscript. We can decompose
\(x\) into two orthogonal parts, with the first part a vector with all
elements equal to the mean and the second part a vector with the
deviations from the mean. Thus \begin{equation}
x_i=x_\bullet+(x_i-x_\bullet).\label{eq-onedimelem}
\end{equation}

Equation \eqref{eq-onedimelem} can be written as \begin{equation}
x=Px+(I-P)x,\label{eq-onedimproj}
\end{equation} where \(P\) projects on the subspace of all constant
vectors, i.e.~\(P:=n^{-1}ee'\) with \(e\) a vector with all elements
equal to one. Now \(P^2=P\) and thus \(P(I-P)=0\), which shows that
\(Px\) and \((I-P)x\) are orthogonal. Also, as Pythagoras used to say,
\begin{equation}
\|x\|^2=\|Px\|^2+\|I-Px\|^2.\label{eq-onedimpyth}
\end{equation}

The projection formulation suggests some obvious generalizations. We can
replace the one-dimensional subspace of all constant vectors by an
arbitrary subspace \(S\), and we can replace the inner product \(x'y\)
by an arbitrary inner product \(x'Ay\), with \(A\) positive
semi-definite. If \(Y\) is an \(A\)-orthonormal basis for the subspace
\(S\), i.e.~\(Y'AY=I\), then \(Px=YY'Ax\).

A further generalization uses multiple mutually orthogonal subspaces
\(S_1,\cdots,S_m\), with corresponding projectors \(P_1,\cdots,P_m\). If
\(P_1+\cdots+P_m=I\) we can write \begin{equation}
x=P_1x+P_2x+\cdots P_mx,
\end{equation} and if \(P_1+\cdots+P_m\lesssim I\) in the Loewner sense
then \begin{equation}
x=P_1x+P_2x+\cdots P_mx+(I-P_1-\cdots-P_m)x
\end{equation} Note that the ``residual'' \(I-P_1-\cdots-P_m\) is
orthogonal to all \(P_1,\cdots,P_m\).

\subsection{Matrices}\label{matrices}

Suppose \(X\) is an \(n\times m\) matrix, an element of the space
\(\mathbb{R}^{n\times m}\) with inner product \(\text{tr}\ X'Y\) and
squared norm \(\|X\|^2=\text{tr}\ X'X\). The decomposition we are
interested in is \begin{equation}
x_{ij}=x_{\bullet\bullet}+(x_{i\bullet}-x_{\bullet\bullet})+(x_{\bullet j}-x_{\bullet\bullet})+
(x_{ij}-x_{i\bullet}-x_{\bullet j}+x_{\bullet\bullet}).\label{eq-twodimelem}
\end{equation} If we define the projectors \begin{subequations}
\begin{align}
P_r&=n^{-1}ee',\\
Q_r&=I-P_r,
\end{align}
and
\begin{align}
P_c&=m^{-1}ee',\\
Q_c&=I-P_c.
\end{align}
\end{subequations} then equation \eqref{eq-twodimelem} becomes
\begin{equation}
X=P_rXP_c+P_rXQ_c+Q_rXP_c+Q_rXQ_c.\label{eq-twodimproj}
\end{equation} Thius clearly shows four projectors, orthogonal to each
other, adding up to the identity.

\subsection{Arrays}\label{arrays}

\begin{multline}
x_{ijk}=x_{\bullet\bullet\bullet}+(x_{i\bullet\bullet}-x_{\bullet\bullet\bullet})+(x_{\bullet j \bullet}-x_{\bullet\bullet\bullet})+(x_{\bullet\bullet k}-x_{\bullet\bullet\bullet})+\\
(x_{ij\bullet}-x_{i\bullet\bullet}-x_{\bullet j\bullet}+x_{\bullet\bullet\bullet})+\\
(x_{i\bullet k}-x_{i\bullet\bullet}-x_{\bullet\bullet k}+x_{\bullet\bullet\bullet})+\\
(x_{\bullet jk}-x_{\bullet j\bullet}-x_{\bullet\bullet k}+x_{\bullet\bullet\bullet})+\\
(x_{ijk}
-x_{ij\bullet}-x_{i\bullet k}-x_{\bullet jk}+x_{i\bullet\bullet}+x_{\bullet j\bullet}+x_{\bullet\bullet k}-x_{\bullet\bullet\bullet})
\end{multline}

\section{Projectors}\label{projectors}

\(\{P_1(A)\}_{ijk}=x_{i\bullet\bullet}\)

\(\{Q_1(A)\}_{ijk}=\delta^{ijk}-x_{i\bullet\bullet}\)

\(\{P_{12}(A)\}_{ijk}=x_{ij\bullet}\)

\section{Least Squares}\label{least-squares}

\section{Interactions}\label{interactions}

\section{Log-linear Analysis}\label{log-linear-analysis}

\section{Orthogonal functions}\label{orthogonal-functions}

Suppose \(M_n\) is the set of matrices \(X\) of order \(n\) with

\begin{itemize}
\tightlist
\item
  first column all elements equal to one, and
\item
  \(X'X=nI\).
\end{itemize}

Suppose the array \(A\) is \(n\times m\times \ell\). Choose
\(X\in M_n\), \(Y\in M_m\), and \(Z\in M_\ell\). Define \[
\alpha_{ijk}=\sum_{i'=1}^n\sum_{j'=1}^m\sum_{k'=1}^l x_{i'j'k'}x_{i'i}y_{j'j}z_{k'k}
\]




\end{document}
